Fix an integer \(n\ge 1\). Let \(f\in\mathbb{C}[z]\) be a monic polynomial of degree \(n\) whose roots \(\alpha_1,\dots,\alpha_n\) satisfy \(|\alpha_j|\le 1\).
Define the lemniscate sublevel set
\[
E(f):=\{z\in\mathbb{C}: |f(z)|\le 1\}.
\]
For such \(f\), define
\[
\ell(f):=\inf\{\operatorname{length}(\gamma): \gamma\text{ is a rectifiable path in }E(f),\ \gamma(0)=0,\ |\gamma(1)|=1\}.
\]
(We interpret ``a path joining \(0\) to \(|z|=1\)'' as a continuous curve with one endpoint at \(0\) and the other on the unit circle; rectifiability is imposed to make ``length'' well-defined.)

The natural extremal quantity is
\[
L_n:=\sup\{\ell(f): f\text{ monic degree }n\text{ with all roots }|\alpha_j|\le 1\}.
\]
The question asks for the behavior of \(L_n\), in particular how large \(\ell(f)\) can be as a function of \(n\).
